% !TEX program = xelatex
% 系统工程理论与实践 LaTeX 空骨架。
% 本文件只提供版式和环境模板，不含实验数据、实验结论或正文论断。
\documentclass{setp-new}

% 此处配置中英文字体。setp-new.cls 使用 fontset=none，须在正文中显式指定字体。
\usepackage{fontspec}
\setmainfont{Times New Roman}
\setsansfont{Arial}
\IfFontExistsTF{Songti SC}{%
  \newcommand{\setpSongFont}{Songti SC}%
}{%
  \newcommand{\setpSongFont}{SimSun}%
}
\IfFontExistsTF{Songti SC}{%
  \newcommand{\setpFangSongFont}{Songti SC}%
}{%
  \newcommand{\setpFangSongFont}{FangSong}%
}
\IfFontExistsTF{SimHei}{%
  \newcommand{\setpHeiFont}{SimHei}%
}{%
  \IfFontExistsTF{Noto Sans CJK SC}{%
    \newcommand{\setpHeiFont}{Noto Sans CJK SC}%
  }{%
    % 若系统没有独立黑体，则以可用宋体作最后回退，保证骨架可编译。
    \newcommand{\setpHeiFont}{\setpSongFont}%
  }
}
\setCJKmainfont[BoldFont=\setpHeiFont]{\setpSongFont}
\setCJKsansfont{\setpHeiFont}
\setCJKmonofont{\setpSongFont}
\setCJKfamilyfont{zhsong}[BoldFont=\setpHeiFont]{\setpSongFont}
\setCJKfamilyfont{zhhei}{\setpHeiFont}
\setCJKfamilyfont{zhfs}{\setpFangSongFont}
\providecommand{\songti}{}
\providecommand{\heiti}{}
\providecommand{\fangsong}{}
\renewcommand{\songti}{\CJKfamily{zhsong}}
\renewcommand{\heiti}{\CJKfamily{zhhei}}
\renewcommand{\fangsong}{\CJKfamily{zhfs}}
\XeTeXlinebreaklocale "zh"
\XeTeXlinebreakskip=0pt plus 1pt

% 此处加载骨架示例所需宏包；setp-new.cls 已加载 amsmath、booktabs 和 graphicx。
\usepackage{mathtools}
\setlength{\parindent}{2em}
\setlength{\headheight}{32pt}
\setlength{\emergencystretch}{3em}

% 此处填写期刊卷期、起始页、DOI、中图分类号和文献标志码；录用前未知项可暂留占位。
% 月份参数必须是 1--12 的整数，不能写中文“待定”，否则类文件无法生成英文月份。
\jolname{系统工程理论与实践}{Systems Engineering --- Theory \& Practice}
\Volume{20XX}{1}{XX}{X}
\PageNum{1}
\PaperDOI{}
\Category{}{A}

% 此处填写中文短题名，用于页眉。
\TitleMark{中文短题名（待填写）}
% 此处填写中文题名。
\Title{中文题名（待填写）}
% 此处填写中文作者姓名及作者单位、城市和邮政编码。
\Author{作者姓名（待填写）}{作者单位、城市和邮政编码（待填写）}
% 此处填写中文摘要；摘要应独立说明研究问题、方法、结果和结论。
\Abstract{中文摘要（待填写）}
% 此处填写中文关键词，关键词之间使用分号。
\Keywords{关键词一；关键词二；关键词三}

% 此处填写英文题名。
\ETitle{English Title (To Be Filled)}
% 此处填写作者姓名拼音或英文写法，以及英文单位和地址。
\EAuthor{AUTHOR NAME (TO BE FILLED)}{Affiliation and address (to be filled)}
% 此处填写英文摘要，与中文摘要信息一致。
\EAbstract{English abstract (to be filled).}
% 此处填写英文关键词，关键词之间使用分号。
\EKeywords{keyword one; keyword two; keyword three}

\begin{document}
\maketitle

\section{引言}

% 此处放研究背景、问题提出、文献回顾、研究不足和本文贡献。
【正文占位：引言】

\section{模型建立}

% 此处放问题描述、基本假设、符号说明、目标函数和约束条件。
【正文占位：模型建立】

% 此处放公式。用 equation 环境自动编号；用 label 与 eqref 在正文中交叉引用。
\begin{equation}
  \min_{x\in\mathcal{X}} f(x)
  \label{eq:template-objective}
\end{equation}
% 此处解释公式中的符号、单位和成立条件。
【正文占位：式~\eqref{eq:template-objective}的符号说明】

% 此处放三线表。表题置于表格上方；单位写入列名，表注放在表格下方。
\begin{table}[htbp]
  \centering
  \caption{三线表标题（待填写）}
  \label{tab:template-three-line}
  \begin{tabular}{lll}
    \toprule
    项目 & 符号或口径 & 数值与单位 \\
    \midrule
    【占位】 & 【占位】 & 【占位】 \\
    \bottomrule
  \end{tabular}
  \tabnote{注：此处填写数据口径、缩写、显著性或其他必要说明。}
\end{table}

\section{算法设计}

% 此处放算法总体框架、关键步骤、复杂度或停止条件。
【正文占位：算法设计】

% 此处放可复现生成的矢量图片；正式文件建议置于 figures/ 目录。
\begin{figure}[htbp]
  \centering
  % 将下一行占位框替换为：\includegraphics[width=0.90\linewidth]{figures/图片文件名.pdf}
  \fbox{\parbox[c][38mm][c]{0.86\linewidth}{\centering 图片占位：在此嵌入矢量 PDF}}
  \caption{图片标题（待填写）}
  \label{fig:template-image}
\end{figure}

\section{数值实验}

% 此处放算例与参数、比较协议、结果表图及与证据一致的分析。
% 未产生并获准进入正文的正式结果不得写入本节。
【正文占位：数值实验】

\section{结语}

% 此处概括经正式结果支持的研究结论、管理启示、局限和后续研究。
【正文占位：结语】

% 此处按《系统工程理论与实践》的著录体例逐条填写参考文献。
% 正文引用示例：\cite{ref:template-cn-journal}。
\begin{thebibliography}{99}
  % 中文期刊论文模板：中文著录之后另起一行给出对应英文著录。
  \bibitem{ref:template-cn-journal}
  中文作者. 中文文献题名[J]. 中文刊名, 年, 卷(期): 起页--止页.\\
  English author. English title[J]. English Journal Title, year, volume(issue): start--end.

  % 外文期刊论文模板：作者. 题名[J]. 刊名, 年, 卷(期): 起页--止页.
  \bibitem{ref:template-en-journal}
  Author. Article title[J]. Journal Title, year, volume(issue): start--end.

  % 专著模板：作者. 书名[M]. 出版地: 出版者, 年.
  \bibitem{ref:template-book}
  作者. 书名[M]. 出版地: 出版者, 年.

  % 网络资源模板：责任者. 题名[EB/OL]. 发布日期[引用日期]. 网址.
  \bibitem{ref:template-online}
  责任者. 网络资源题名[EB/OL]. 发布日期[引用日期]. \url{https://example.invalid}.
\end{thebibliography}

\end{document}
